% ---------------------------------------------------------------------------
% Carta de motivación — Boot Camp Legislativo 2026 (Asamblea Nacional)
% Manuel Josué Malla Campoverde
% Mismo membrete que cv/cv-bootcamp-legislativo-2026-es.tex
% Texto fuente: carta-motivacion.md
%
% Compilar con:  pdflatex carta-motivacion.tex
% ---------------------------------------------------------------------------
\documentclass[letterpaper,11pt]{article}

\usepackage[utf8]{inputenc}
\usepackage[T1]{fontenc}
\usepackage{lmodern}
\usepackage[spanish,es-noshorthands,es-nolists,es-noquoting]{babel}
\usepackage[left=0.9in,right=0.9in,top=0.6in,bottom=0.6in]{geometry}
\usepackage[hidelinks]{hyperref}
\usepackage{setspace}

\pagestyle{empty}
\setlength{\parindent}{0pt}
\setlength{\parskip}{5pt}

\begin{document}

% ---------------------------- Membrete -------------------------------------
\begin{center}
  {\large \textbf{MANUEL JOSUÉ MALLA CAMPOVERDE}}\\[2pt]
  {\small Machala, El Oro, Ecuador $\diamond$
   \href{mailto:mmalla1@outlook.com}{mmalla1@outlook.com} $\diamond$
   +593 97 948 9503 $\diamond$
   \href{https://jossuema.tech}{jossuema.tech}}
\end{center}
\vspace{-4pt}
\rule{\textwidth}{0.9pt}
\vspace{2pt}

% ---------------------------- Encabezado -----------------------------------
Machala, 6 de septiembre de 2026

Señores\\
\textbf{Comité de Selección --- Boot Camp Legislativo 2026}\\
Asamblea Nacional del Ecuador\\
Presente.

De mi consideración:

% ------------------------------ Cuerpo -------------------------------------
Me dirijo a ustedes para expresar mi interés en participar en el Boot Camp
Legislativo 2026. Soy Manuel Josué Malla Campoverde, estudiante de Ingeniería
en Tecnologías de la Información en la Universidad Técnica de Machala y
representante estudiantil ante la Asamblea de Educación Superior de mi
universidad.

Creo firmemente que la tecnología puede transformar realidades, y durante mi
formación he buscado llevar esa convicción a la práctica. Desarrollé una
investigación sobre la distribución territorial de los homicidios en el
Ecuador a partir de datos gubernamentales, con el fin de identificar las zonas
donde el fenómeno se concentra con mayor intensidad. Participé en el
desarrollo del Buzón Inteligente de mi universidad, plataforma que recoge
quejas, sugerencias y propuestas de la comunidad universitaria. Y dicté un
curso introductorio de programación para jóvenes en alianza con el Municipio
de Huaquillas. Esas experiencias me enseñaron dos cosas: que los problemas
públicos se comprenden mejor cuando se los mira con evidencia, y que la
tecnología sirve de verdad cuando llega a las personas.

También he participado en espacios de liderazgo y representación estudiantil,
porque entendí que generar impacto no consiste solamente en crear soluciones,
sino en involucrar a más personas en ellas. Como representante estudiantil
gestiono alianzas con instituciones públicas y privadas para respaldar
iniciativas educativas, y ahí he visto que las buenas ideas
necesitan además reglas, acuerdos y procesos que las sostengan.
Ese es el vacío que quiero llenar en mi formación.

Del programa espero aprender sobre participación ciudadana, liderazgo y
procesos legislativos: cómo se construye una norma, de qué manera la
ciudadanía puede incidir formalmente en esas decisiones y cómo se traduce una
necesidad social en una propuesta viable. Quiero comprender
cómo transformar ideas en acciones que contribuyan al desarrollo de nuestro
país.

Me comprometo a replicar lo aprendido en mi ámbito de acción. En primer lugar,
en la Asamblea de Educación Superior de mi universidad,
llevando mecanismos de participación mejor estructurados al espacio donde ya
represento a mis compañeros. En segundo lugar, mediante talleres para
estudiantes, como los que facilito en la comunidad de Embajadores Estudiantiles
de Microsoft Learn, esta vez dedicados a la participación ciudadana y al uso de
datos públicos. Y en tercer lugar, extendiendo esa formación a jóvenes de la
provincia de El Oro, como ya lo hice en Huaquillas junto al gobierno local.

Mi objetivo es que más jóvenes en el Ecuador se sumen a la transformación del
país, utilizando la tecnología como una herramienta para generar oportunidades
y bienestar. Agradezco la oportunidad de ser considerado para este programa.

% ------------------------------ Cierre -------------------------------------
\vspace{6pt}
Atentamente,

\vspace{20pt}
\textbf{Manuel Josué Malla Campoverde}\\
{\small Estudiante de Ingeniería en Tecnologías de la Información\\
Universidad Técnica de Machala\\
mmalla1@outlook.com $\cdot$ +593 97 948 9503}

\end{document}
